\chapter{Protein-Interaction Prediction: Data, Features, and Predictive Pipeline} \label{chap:protein_pipeline}
\chaptersummary{The protein application begins by defining what counts as a positive interaction in this study and by building the surrogate used later for inference. This chapter describes the benchmark splits, the representation choices, and the predictive pipeline that turns protein pairs into a fixed edge-level score as predictive outcomes.}


In this chapter, we investigate a simple and concrete biological question: Given a pair of proteins, we want to predict whether the benchmark classifies that pair as a physical protein--protein interaction. Positive edges therefore correspond to curated physical interactions, whereas negative edges are defined either by benchmark (Intra datset) or by sampled non-positive pairs (BIOGRID human dataset). The point here is not to catalog implementation details. The goal of this section is not to enumerate implementation details, but to show how this biological task is formulated as a leakage-aware prediction problem whose outputs can later be used in the downstream inferential analysis.

The predictive pipeline follows a simple sequence. Protein sequences and UniProt text descriptions are encoded into protein-level embeddings, those embeddings are converted into pairwise edge features, several prediction heads are compared, and a fixed three-seed attention ensemble is taken forward as the default surrogate. Figure~\ref{fig:chap4_pipeline_overview} summarizes that workflow. The discussion follows the same order: train/test split design and benchmark construction first, protein representations and pairvise features second, and lastly the final surrogate outcomes.

\begin{figure}[htbp]
\centering
\MaybeIncludeGraphics[width=\textwidth]{Pipeline_thesis.pdf}
\caption[Protein-edge predictive pipeline]{High-level predictive pipeline used in the protein application. Protein sequences and UniProt descriptions are encoded into protein-level embeddings, converted into edge-level pair features, and passed to baseline and neural prediction heads. The final surrogate used in later chapters is the frozen three-seed attention ensemble, whose averaged edge-level probability is exported as the prediction input for downstream prediction-powered inference and \texttt{PPI++} analysis.}
\label{fig:chap4_pipeline_overview}
\end{figure}

\section[Data Sources and Splits]{Data Sources, Curation, and Split Construction}

Two datasets anchor the protein application.
\begin{enumerate}
    \item \textbf{Intra}: the internal benchmark used for model selection and the main in-domain evaluation.
    \item \textbf{BIOGRID human physical interactions}: an external benchmark used to test robustness under a protein-disjoint split.
\end{enumerate}

The choice reflects the broader literature on interactome construction. BioGRID is one of the major curated resources for experimentally observed interactions, whereas the human reference interactome and related mapping projects highlight both the value and the incompleteness of current protein-protein interaction catalogs \parencite{biogrid2019, luck2020, vidal2011}. In this study, the raw edge tables are first normalized into processed pair datasets and then split into train, validation, and test partitions with explicit leakage controls.

One difference between the two datasets should be stated explicitly because it affects how the benchmark is interpreted. The Intra benchmark already provides separate positive and negative edge files, so the negative pairs are predefined by the benchmark itself. The BIOGRID-human dataset, by contrast, contributes curated positive physical interactions only; its negative edges are therefore constructed synthetically by sampling protein pairs from the mapped BIOGRID protein universe whose pair does not appear in the curated positive set. These BIOGRID negatives should be understood as sampled non-positive pairs, not as experimentally validated non-interactions.

\begin{table}[ht]
\centering
\small
\renewcommand{\arraystretch}{1.12}
\setlength{\tabcolsep}{5pt}
\begin{tabular}{@{}llrrrl@{}}
\toprule
Dataset & Split & Pos. & Neg. & Total & Role \\
\midrule
\multirow{3}{*}{Intra}
  & Train & 81{,}596 & 81{,}596 & 163{,}192 & Train \\
  & Val.  & 29{,}630 & 29{,}630 & 59{,}260 & Pilot / val. \\
  & Test  & 26{,}024 & 26{,}024 & 52{,}048 & Main test \\
\addlinespace[2pt]
\multirow{3}{*}{BIOGRID}
  & Train & 394{,}829 & 399{,}610 & 794{,}439 & Pool cand. \\
  & Val.  & 19{,}814 & 18{,}160 & 37{,}974 & Ext. val. \\
  & Test  & 17{,}819 & 18{,}950 & 36{,}769 & Ext. test \\
\addlinespace[2pt]
Merged pool & Train & 416{,}686 & 480{,}172 & 896{,}858 & Merged train \\
\bottomrule
\end{tabular}
\caption[Processed data splits for the protein application]{Summary of the processed splits used in the protein application. The merged training pool is formed after removing overlaps with validation and test holdouts and resolving duplicate-label conflicts with a positive-wins rule.}
\label{tab:chap4_data_splits}
\end{table}

Table~\ref{tab:chap4_data_splits} is the central reproducibility summary for the application. The Intra validation split serves two roles later in the thesis: it is the model-selection split for predictive work and the pilot/calibration split for several planning calculations. The BIOGRID validation and test splits are held out from the merged training pool and are constructed to reduce leakage through shared proteins. What matters for the thesis is the split design itself rather than the particular pre-processing scripts used to realize it.

\section[Representations and Pair Features]{Per-Protein Representations and Pair Features}

Each protein is represented from two complementary views.
\begin{enumerate}
    \item \textbf{Sequence view}: embeddings from pretrained protein language models, including ProtBERT and ESM2 model variants.
    \item \textbf{Text view}: embeddings from UniProt description text using a compact sentence-transformer backbone.
\end{enumerate}


The sequence representation is motivated by the strong transfer performance of large self-supervised protein language models. ProtTrans showed that large-scale protein pretraining yields broadly useful representations, while the ESM line demonstrated that these models can encode rich structural and functional information \parencite{prottrans2022, lin2023}. The text representation is deliberately complementary. UniProt descriptions contain curated biological context that is not present in raw sequence alone, and sentence-embedding models derived from MiniLM provide a compact way to encode that context \parencite{uniprot2025, minilm2020, reimers-2021-all-minilm-l6-v2}.
The predictive unit is an edge, so protein-level embeddings must be converted into pair-level features. For a protein pair $(u,v)$, the workflow uses the operator-based representation
\[
[u+v,\; |u-v|,\; u \odot v,\; \cos(u,v)],
\]
applied to the sequence block, the text block, or the concatenated sequence-plus-text block. Each operator contributes a different summary of the pair. The sum \(u+v\) captures pooled shared signal or joint magnitude, \(|u-v|\) captures elementwise mismatch, \(u \odot v\) captures coordinate-wise interaction or alignment, and \(\cos(u,v)\) adds one global directional-similarity score. The first three operators are vector-valued feature blocks, whereas the cosine term is appended as a scalar summary. This choice is pragmatic. It is simple enough to support controlled ablations and downstream uncertainty analysis, while still giving the neural heads access to nonlinear interactions among complementary summary operators. The implementation realizes this construction explicitly after the embedding stage, but the statistical object that matters here is the resulting pair-level surrogate input.

In Figure~\ref{fig:chap4_pipeline_overview}, this is the step labeled feature construction. The important point is that the unit has changed: the raw scientific objects are individual proteins, but the predictive target is an edge. The feature map therefore has to convert two protein-level representations into one edge-level surrogate input.

\section[Protein-Interaction Prediction Literature]{Relationship to the Protein-Interaction Prediction Literature}

Chapter~\ref{chap:introduction} has already summarized the broader shift in the protein–protein interaction literature from sequence-based prediction toward protein language model representations and hybrid architectures. The aim of the present section is narrower. Literature shows that strong single-protein representations are now feasible, but it also shows that pair-input evaluation is highly sensitive to split construction and negative-edge design \parencite{sun2017, dscript2021, xcapt52024, lee2023, park2012, bernett2024}. The current thesis does not claim a universal state of the art against those studies. Instead, it takes from them a more defensible methodological lesson: representation quality matters, but benchmark design matters just as much. For that reason, the analysis here prioritizes leakage-aware splits, external robustness checks, and prepare for downstream statistical inference using the PPI framework mentioned in Chapter~\ref{chap:foundation_llm}.

\section[Training Pool and Neural Predictors]{Merged Training Pool and Neural Predictors}

The mainline training set is not any single raw split. It is a merged, leakage-controlled training pool formed from the train partitions of the processed Intra and BIOGRID datasets after removing overlaps with all downstream holdouts. In the present analysis, this yields 896{,}858 training edges, of which 416{,}686 are positive and 480{,}172 are negative. The pool-construction step also removes 26{,}830 pairs that overlap with later validation or test sets and resolves 594 duplicate-label conflicts with a positive-win rule: if a protein pair appears with both labels after merging sources, the merged pool retains it as positive, because curated interaction catalogs are incomplete and sampled non-positive pairs should not be interpreted as experimentally verified non-interactions \parencite{luck2020, park2012, bernett2024}. These details matter because the later inference claims depend on the surrogate being trained without contaminating the holdout targets.

Three neural heads are compared on top of the pair features.
\begin{enumerate}
    \item \textbf{NN-LogReg}: a linear head that serves as the simplest neural baseline.
    \item \textbf{NN-MLP}: a multilayer perceptron that allows nonlinear mixing of the pair features.
    \item \textbf{NN-Attn}: a multihead self-attention head operating over operator/modality tokens.
\end{enumerate}

For the \texttt{seq\_text} setting, the attention model does not treat the full feature vector as one undifferentiated block. Instead, it groups the pair representation into sequence and text operator tokens, projects those tokens into a shared latent space, applies multihead self-attention, and then produces one edge-level probability. Figure~\ref{fig:chap4_pipeline_overview} shows these three heads as parallel model families, with the attention branch feeding the final ensemble. This is a deliberately moderate architecture choice relative to full end-to-end cross-encoders: the thesis needs a strong and stable surrogate that can be exported cleanly into downstream inference, not only a maximally complex predictor.

\section[Frozen Model and Handoff]{Frozen Best Model and Prediction Handoff}

The final surrogate carried into the later chapters is
\begin{center}
\texttt{seq\_text + esm2\_3b + attention ensemble}.
\end{center}
It is built by training three attention models with different random seeds and averaging their predicted probabilities. On the held-out test splits, the ensemble achieves PR-AUC 0.715 on Intra and PR-AUC 0.799 on BIOGRID-human, making it the strongest practical surrogate used in the later inference layer.

The predictive and inferential workflows are connected through an explicit export contract. Training produces seed-level predictions, the ensemble step averages those probabilities, and a canonical export step writes the resulting edge-level probabilities to a stable file schema for the inference layer. In this thesis, ``frozen'' means that model selection has ended and that this exported prediction object is treated as fixed input thereafter. This separation is important because it keeps the prediction problem and the inference problem distinct: the neural model supplies the surrogate score, while prediction-powered procedures use a smaller labeled subset to correct and quantify uncertainty.

\section{Chapter Summary}

The protein application therefore rests on four design decisions: leakage-aware split construction, complementary sequence and text representations, operator-based pair features, and a frozen attention ensemble used as the default surrogate. These choices are informed by curated protein-protein interaction resources, protein language modeling, and the recent literature on leakage-aware sequence-based protein-interaction prediction \parencite{biogrid2019, luck2020, prottrans2022, dscript2021, park2012, bernett2024}. The next chapter evaluates that surrogate as a prediction layer and then uses its exported scores for prediction-powered inference and \texttt{PPI++} estimation, labeled-budget planning, and dependence sensitivity analysis.
